\documentclass[11pt,a4paper]{article}
\usepackage[T1]{fontenc}
\usepackage[utf8]{inputenc}
\usepackage{lmodern}
\usepackage{microtype}
\usepackage[a4paper,margin=27mm,headheight=14pt]{geometry}
\usepackage{amsmath,amssymb,amsthm,mathtools}
\usepackage{booktabs,array,longtable}
\usepackage{xcolor}
\usepackage{enumitem}
\usepackage{fancyhdr}
\usepackage{hyperref}
\usepackage{bookmark}
\hypersetup{colorlinks=true,linkcolor=blue!45!black,citecolor=blue!45!black,
 urlcolor=blue!45!black,pdftitle={Arithmetic Convolution Factors of Fabius-Type Laws},
 pdfauthor={Research manuscript prepared for the ProveIt programme},
 pdfsubject={Divisibility ladders, singular remainders, Borel classification and computability}}
\urlstyle{same}
\pagestyle{fancy}\fancyhf{}
\fancyhead[L]{\small Arithmetic convolution factors}
\fancyhead[R]{\small ProveIt research manuscript}
\fancyfoot[C]{\thepage}
\setlength{\parindent}{1.3em}
\setlength{\parskip}{3pt}
\setlist{itemsep=3pt,topsep=5pt}
\allowdisplaybreaks[2]
\emergencystretch=2em
\newtheorem{theorem}{Theorem}[section]
\newtheorem{lemma}[theorem]{Lemma}
\newtheorem{proposition}[theorem]{Proposition}
\newtheorem{corollary}[theorem]{Corollary}
\theoremstyle{definition}
\newtheorem{definition}[theorem]{Definition}
\newtheorem{example}[theorem]{Example}
\newtheorem{question}{Research question}
\theoremstyle{remark}
\newtheorem{remark}[theorem]{Remark}
% ed. (2026-09-29): unnumbered environment for editorial notes.
\newtheorem*{ednote}{Editorial note (ProveIt, 2026-09-29)}
% ed. (2026-09-30): a second unnumbered environment for the notes of 2026-09-30.
\newtheorem*{ednotelater}{Editorial note (ProveIt, 2026-09-30)}
\newcommand{\R}{\mathbb R}
\newcommand{\N}{\mathbb N}
\newcommand{\Z}{\mathbb Z}
\newcommand{\C}{\mathbb C}
\newcommand{\E}{\mathbb E}
\newcommand{\Prob}{\mathbb P}
\newcommand{\Law}{\operatorname{Law}}
\newcommand{\supp}{\operatorname{supp}}
\newcommand{\sinc}{\operatorname{sinc}}
\newcommand{\ord}{\operatorname{ord}}
\newcommand{\lcm}{\operatorname{lcm}}
\newcommand{\Var}{\operatorname{Var}}
\newcommand{\BV}{\operatorname{BV}}
\newcommand{\D}{\mathsf D}
\newcommand{\Unif}{\mathsf U}
\newcommand{\Mu}[1]{\mu^{[#1]}}
\newcommand{\etaam}{\eta_{a,m}}
\newcommand{\sourcecommit}{\texttt{fbba58593dc0622aa914972896150d4848f935b5}}

\begin{document}
% ed. (2026-09-29): no page anchor on the title page (it duplicated page.1).
\hypersetup{pageanchor=false}
\begin{titlepage}
\vspace*{12mm}
\begin{center}
{\small RESEARCH ARTICLE \quad\textbar\quad 28 SEPTEMBER 2026}\par
\vspace{13mm}
{\LARGE\bfseries Arithmetic Convolution Factors\\[5pt]
of Fabius-Type Laws}\par
\vspace{8mm}
{\large Exact divisibility, fractal remainders,\\[3pt]
and classification barriers}\par
\vspace{12mm}
{\normalsize Prepared for the ProveIt research programme}
\end{center}
\vspace{9mm}
\begin{abstract}
Let $A_0=1$ and let $A_k$ be a strictly increasing chain of positive
integers under divisibility. The random series
$X_A=\sum_{k\ge0}U_k/A_k$, with independent centered unit uniforms $U_k$,
has a compactly supported $C^\infty$ density. We give an exact classification
of factorizations $\mu_A=\D_c\mu_B*\nu$ by probability measures:
they exist precisely when $c=1/m$ for a positive integer $m$ and
$A_k\mid mB_k$ for every $k$. The unique remainder has an explicit
independent-digit representation. This converts a positive-definiteness
problem for entire Fourier quotients into coordinatewise arithmetic.

The classification yields a rigid cross-base theorem for geometric-uniform
laws, an embedding of inclusion modulo finite sets into scaled convolution
factorability, and a compact family of smooth convolution divisors of the single
Rvachev law on which mutual factorability has no Borel real-valued complete
invariant. With
primitive-recursive ladder presentations, fixed-scale factorability is
$\Pi^0_1$-complete and existence of some scale is $\Sigma^0_2$-complete.
Explicit spline approximation bounds show that these barriers persist for
uniformly computable smooth densities. For same-base geometric laws we
also classify the remainder as singular continuous, a finite uniform
spline, or a finite smoothing of a singular measure; we prove sharp
integer regularity, sharp power Fourier decay, and a sharp H\"older result
in the separated-digit case. Further questions isolate the boundaries
of the classification rather than presuming that they are already solved.
\end{abstract}
\vfill
\noindent\textbf{Status and scope.}
All asserted mathematical results below have written proofs. The auxiliary
program performs exact finite regression checks, not proof-assistant
verification. No Lean verification or peer review is claimed. Classical
inputs and identified repository overlaps are distinguished from the
extensions developed here. Priority of the latter has not been established
by an exhaustive literature or repository audit.
\end{titlepage}
% ed. (2026-09-29): page anchors resume after the title page.
\hypersetup{pageanchor=true}

\begingroup
\small
\setlength{\parskip}{0pt}
\tableofcontents
\endgroup
\newpage

\section{Research question, provenance, and main conclusions}
\label{sec:introduction}

\subsection{The question behind the Fourier product}
The Rvachev up-function is a particularly structured compactly supported
smooth probability density. Its Fourier transform is an infinite product
of sinc functions. In the normalization used here,
\begin{equation}\label{eq:up-product}
 \widehat{\mathrm{up}}(t)=\prod_{k=0}^{\infty}
 \frac{\sin(\pi t/2^k)}{\pi t/2^k}.
\end{equation}
Arias de Reyna gives this product and the integer-zero multiplicity formula
$1+v_2(n)$ explicitly \cite{arias}. These formulas, and the probabilistic
construction underlying them, are established starting points, not new
claims of this article. Fabius's original construction provides the
historical probabilistic context \cite{fabius}.

The product suggests a question that is more restrictive than formal
power-series division. Given two such laws, when does removing a scaled
copy of one from the other leave a \emph{probability measure}? Precisely,
for which $c>0$ is there a probability $\nu$ such that
\begin{equation}\label{eq:factor-question}
 \mu=\D_c\rho*\nu?
\end{equation}
Here $\D_c\rho$ is the law of $cY$ when $Y$ has law $\rho$, and $*$ is
convolution, equivalently addition of independent random variables.
A formally defined Fourier quotient need not be a characteristic function.
Even an entire quotient is not automatically positive definite. The
central result is that for a natural class of nested integer denominators,
these analytic and positivity questions have a complete arithmetic answer.

\subsection{Connection with the repository and limits of the audit}
The repository was inspected at commit \sourcecommit.
Its Fabius research documentation discusses geometric-uniform laws,
Fourier and Laplace products, moment reciprocity, and deconvolution
\cite{repo-inverse}. Its Shape, Divisibility, and Stein Geometry package
already includes identical-convolution rootlessness for the Fabius--Rvachev
law \cite{repo-shape}. We therefore do not present the dyadic zero formula,
probabilistic construction, or dyadic rootlessness as repository-new results.

Our principal extension is to two arbitrary divisibility ladders, and then
to the order-theoretic and effective complexity encoded by those ladders.
The same-base remainder analysis supplies a complementary analytic
classification. The general geometric self-factor theorem is included as
a useful bridge and is proved independently; no priority claim is attached
to that elementary product consequence.

The audit was scoped, not exhaustive. In particular, direct retrieval of
the large geometric-$q$ frontier TeX file exceeded the connector's supported
response size. Navigation records and smaller relevant files were inspected,
but absence of a theorem from those records is not evidence that it occurs
nowhere in the larger corpus. The article thus establishes mathematical
statements and gives their proofs, without claiming to settle a named
long-standing conjecture whose open status has not been verified. The
classical descriptive-set-theoretic obstruction used later is also not new:
the contribution there is its explicit realization by these smooth laws,
not the nonsmoothness of finite-change equivalence itself \cite{hkl}.

\subsection{What is resolved here}
The first classification is the following exact criterion. Suppose
$A_0=B_0=1$, $A_k\mid A_{k+1}$ and $B_k\mid B_{k+1}$, with both
chains strictly increasing. Then
\begin{equation}\label{eq:main-preview}
 \mu_A=\D_c\mu_B*\nu
 \quad\Longleftrightarrow\quad
 \left\{\begin{array}{l}
 c=1/m\text{ for an integer }m\ge1,\\
 A_k\mid mB_k\text{ for every }k\ge0.
 \end{array}\right.
\end{equation}
This is Theorem~\ref{thm:ladder}. The implication from left to right uses
zero multiplicities; the reverse implication constructs positive digit
measures and therefore resolves, rather than assumes, positive definiteness.

Several consequences make the criterion more than an isolated product
identity. An integer-base target can contain a geometric-uniform source
only when the reciprocal source ratio is an integer multiple of the target
base (Theorem~\ref{thm:crossbase}). A single compact family of smooth densities
realizes inclusion modulo finite sets (Theorem~\ref{thm:boolean}) and
finite-change equivalence (Theorem~\ref{thm:nonsmooth}). Exact effective
classification has specified arithmetical-hierarchy complexity even though
the densities themselves are uniformly approximable with every derivative
(Theorem~\ref{thm:complexity}). Finally, the explicit remainders exhibit a
sharp division between singular measures and finite levels of smoothing
(Theorem~\ref{thm:regularity}).

The results concern genuine independent convolution factors. They do not
classify signed deconvolution kernels, arbitrary nonlinear transports,
approximate factorizations, or every compactly supported smooth law.

\section{Probability laws, entire products, and finite approximations}
\label{sec:preliminaries}

\subsection{Definitions and normalization}
Throughout, $\N_0=\{0,1,2,\ldots\}$ and $\N_{>0}=\{1,2,\ldots\}$.
Let $\Unif_L$ denote uniform probability on $[-L/2,L/2]$, for $L>0$.
Let $(U_k)_{k\ge0}$ be independent random variables with law $\Unif_1$.
Our Fourier convention and normalized sinc are
\[
 \widehat\sigma(t)=\int_\R e^{2\pi itx}\,d\sigma(x),\qquad
 \sinc z=\begin{cases}\sin(\pi z)/(\pi z),&z\ne0,\\1,&z=0.\end{cases}
\]
All laws considered as targets are symmetric, so the opposite Fourier sign
gives the same products.

\begin{definition}[Divisibility ladder]
A divisibility ladder is a sequence $A=(A_k)_{k\ge0}$ of positive integers
such that $A_0=1$, $A_k\mid A_{k+1}$, and $A_{k+1}>A_k$. Set
\begin{equation}\label{eq:ladder-law}
 X_A=\sum_{k\ge0}\frac{U_k}{A_k},\qquad
 \mu_A=\Law(X_A),\qquad R_A=\frac12\sum_{k\ge0}A_k^{-1}.
\end{equation}
For a real ratio $0<q<1$, separately set
\begin{equation}\label{eq:geometric-law}
 X_q=\sum_{k\ge0}q^kU_k,\qquad \mu_q=\Law(X_q).
\end{equation}
For an integer $a\ge2$, abbreviate $\Mu a=\mu_{1/a}$; this is the ladder
law for $A_k=a^k$. In particular, $\Mu2$ has density $\mathrm{up}$ in
\eqref{eq:up-product}, and is not its cumulative distribution function.
\end{definition}

The ladder ratios are integers at least two, so $A_k\ge2^k$.
Both random series converge absolutely for every choice of the bounded
coordinates. In the ladder case $R_A\le1$; in the geometric case the
support radius is $1/(2(1-q))$.

\begin{lemma}[Transforms, support, and smoothness]\label{lem:product}
For a divisibility ladder,
\begin{equation}\label{eq:ladder-product}
 \Phi_A(z):=\widehat\mu_A(z)=\prod_{k\ge0}\sinc(z/A_k)
\end{equation}
converges locally uniformly on $\C$, defines an entire function, and has no
zeros other than the zeros of its individual factors. The law $\mu_A$ has
a $C^\infty$ density $f_A$ with support $[-R_A,R_A]$.
The same assertions hold for $\mu_q$, with factors $\sinc(q^kz)$ and
support $[-1/(2(1-q)),1/(2(1-q))]$.
\end{lemma}
\begin{proof}
Uniform convergence of the random series implies convergence of its
characteristic functions by bounded convergence. On a compact disk,
$\sinc(z/A_k)=1+O(A_k^{-2})$ uniformly, and $\sum A_k^{-2}<\infty$.
The usual product criterion follows directly by taking the analytic
logarithm of sufficiently late factors: their deviations from one have
summable suprema. Thus the tail product is analytic and nonzero on each
small disk avoiding zeros of the finitely many initial factors. This proves
both local uniform convergence and the assertion about zeros.

For every positive integer $M$ and real $|t|\ge1$,
\begin{equation}\label{eq:rapid-decay}
 |\Phi_A(t)|\le \prod_{k=0}^{M-1}\frac{A_k}{\pi|t|}
 =\frac{\prod_{k<M}A_k}{\pi^M|t|^M},
\end{equation}
since the other factors have absolute value at most one. Consequently
$t^r\Phi_A(t)$ is integrable for every $r\ge0$. Fourier inversion and
termwise differentiation under the integral give a smooth density.

The series is confined to $[-R_A,R_A]$. For completeness its support is
the entire interval. The possible sums of the first $N+1$ support intervals
form $[-R_{A,N},R_{A,N}]$, where $R_{A,N}\uparrow R_A$.
The tail has zero in its support: its first finitely many variables can be
made arbitrarily small with positive probability, and its remaining
absolute bound tends to zero. Any interior point of $[-R_A,R_A]$ therefore
lies in the support of a sufficiently long prefix plus its independent
tail. Taking closures includes the endpoints. Replace $A_k^{-1}$ by $q^k$
for the geometric case; all arguments are unchanged.
\end{proof}

\begin{lemma}[Zero multiplicities]\label{lem:zeros}
The zero set of $\Phi_A$ is $\Z\setminus\{0\}$, and for $n\ne0$,
\begin{equation}\label{eq:zero-orders}
 \ord_{z=n}\Phi_A(z)=v_A(n):=\#\{k\ge0:A_k\mid n\}.
\end{equation}
The zero at $1$ is simple. For a geometric law, the zero set is
$\{nq^{-k}:n\in\Z\setminus\{0\},\ k\ge0\}$, and the zero at $1$
is again simple.
\end{lemma}
\begin{proof}
The $k$th ladder factor vanishes simply exactly at nonzero multiples of
$A_k$. Only finitely many factors vanish at a fixed point. Their orders
add, and the remaining product is nonzero by Lemma~\ref{lem:product}.
The first factor accounts for every nonzero integer, and later factors
produce no noninteger zeros. Only the first factor vanishes at $1$.
The same argument applies to the geometric factors; $0<q^k<1$ for $k>0$
prevents any additional vanishing at $1$.
\end{proof}

\subsection{Why compactness makes zero multiplicities legitimate}
\begin{lemma}[Compact factor and cancellation]\label{lem:compact}
Suppose compactly supported probability measures $\mu,\rho$ satisfy
$\mu=\rho*\nu$ for a probability measure $\nu$. Then $\nu$ is compactly
supported, the Fourier identity extends to entire functions, and
\[
 \ord_z\widehat\mu\ \ge\ \ord_z\widehat\rho
\]
at every zero of $\widehat\rho$. If $\widehat\rho$ is nonzero except at
isolated real points, the factor $\nu$ is unique. If $\mu$ and $\rho$
are symmetric, this unique factor is symmetric.
\end{lemma}
\begin{proof}
For support points $x\in\supp\rho$ and $y\in\supp\nu$, every neighborhood
of $x+y$ contains the sum of two positive-measure neighborhoods; independence
therefore gives $x+y\in\supp\mu$. Fixing one $x$ bounds every possible
$y$, so $\nu$ has bounded closed support. The transforms are entire and
agree on the real axis, hence everywhere. Orders of zeros add under
multiplication.

For uniqueness, divide the characteristic functions where
$\widehat\rho(t)\ne0$. The resulting equality for two candidate factors
extends to every real $t$ by continuity, and Fourier uniqueness gives equality
of measures. Reflection of a factor is another factor when both target and
source are symmetric, proving symmetry by uniqueness. The same support
argument, applied to minimum and maximum points of compact supports, shows
that the extreme support points add under convolution.
\end{proof}

\begin{corollary}[Identical-convolution rootlessness]\label{cor:rootless}
None of the laws $\mu_A$ or $\mu_q$ is $\rho^{*j}$ for a probability
measure $\rho$ and an integer $j\ge2$.
\end{corollary}
\begin{proof}
A putative root is compact by the support argument. Its entire transform
raised to the $j$th power would have a zero at $1$ of order divisible by
$j$, contradicting Lemma~\ref{lem:zeros}.
\end{proof}
\noindent The dyadic instance of this short obstruction is already covered
by the repository's divisibility report \cite{repo-shape}. We include the
uniform extension only to delimit the meaning of factorization here:
nonidentical factors can exist even though identical roots do not.

\subsection{Explicit approximation in every derivative}
\begin{proposition}[Effective spline approximation]\label{prop:approximation}
Let $f_{A,N}$ be the density of $\sum_{k=0}^NU_k/A_k$. For $r\ge0$ and
$N\ge r+2$, put
\[
 L_r(A)=2^{r+1}\prod_{j=0}^{r+1}A_j.
\]
Then
\begin{equation}\label{eq:approximation}
 \|f_A^{(r)}-f_{A,N}^{(r)}\|_\infty
 \le \frac{L_r(A)}{2A_N}.
\end{equation}
Every computable ladder therefore yields a computable $C^\infty$ density,
with a modulus supplied explicitly for each derivative.
\end{proposition}
\begin{proof}
Let $g$ be the density of the first $r+3$ uniforms, indexed $0$ through
$r+2$. The distributional derivative of the density of $\Unif_{1/A_j}$
is $A_j(\delta_{-1/(2A_j)}-\delta_{1/(2A_j)})$, with total variation
$2A_j$. Differentiate the factors $j=0,\ldots,r$, leaving two uniform
factors undifferentiated. Their convolution is continuous and has sup norm
at most $A_{r+1}$. It follows that $g^{(r+1)}$ is continuous with norm at
most $L_r(A)$, and $g^{(r)}$ is globally $L_r(A)$-Lipschitz.
Convolution with any remaining prefix preserves this Lipschitz bound.

The tail $T_N=\sum_{k>N}U_k/A_k$ obeys
\[
 |T_N|\le\frac12\sum_{h\ge1}\frac1{A_{N+h}}
 \le\frac1{2A_N}.
\]
Since $f_A^{(r)}(x)=\E f_{A,N}^{(r)}(x-T_N)$, the Lipschitz estimate gives
\eqref{eq:approximation}.

To make computability explicit, write $\ell_j=1/A_j$ and
$S_N=\sum_{j=0}^N\ell_j$. The prefix has the finite spline formula
\begin{equation}\label{eq:box-spline}
 f_{A,N}(x)=\frac1{N!\prod_{j=0}^N\ell_j}
 \sum_{E\subseteq\{0,\ldots,N\}}(-1)^{|E|}
 \left(x+\frac{S_N}{2}-\sum_{j\in E}\ell_j\right)_+^{N}.
\end{equation}
This follows by repeated convolution of interval indicators, or by applying
the corresponding finite differences to $x_+^N/N!$. It is a rational
piecewise polynomial with computable knots. Choose $N$ until the right side
of \eqref{eq:approximation} is smaller than the requested error; this search
terminates because $A_N\ge2^N$. Differentiating the rational spline gives
uniform approximants for each prescribed derivative. No efficiency claim
is made for the exponential-size subset formula.
\end{proof}

\section{A complete classification of ladder factors}
\label{sec:ladder}

\subsection{The positive splitting identity}
\begin{lemma}[One uniform splits into a smaller uniform and digits]
\label{lem:uniform-split}
For an integer $M\ge1$ and $L>0$, let $\ell=L/M$. Then
\begin{equation}\label{eq:uniform-split}
 \Unif_L=\Unif_\ell *
 \left(\frac1M\sum_{j=0}^{M-1}
 \delta_{\ell(j-(M-1)/2)}\right).
\end{equation}
\end{lemma}
\begin{proof}
The translated intervals of length $\ell$ are adjacent and tile
$[-L/2,L/2]$, up to endpoints of Lebesgue measure zero. On each interval
the mixture density is $(1/M)(1/\ell)=1/L$. This proves the probability
identity directly, including $M=1$ when the digit law is $\delta_0$.
\end{proof}

\subsection{The main theorem}
\begin{theorem}[Arithmetic convolution classification]\label{thm:ladder}
Let $A$ and $B$ be divisibility ladders and let $c>0$. There is a probability
measure $\nu$ satisfying
\[
 \mu_A=\D_c\mu_B*\nu
\]
if and only if there is an integer $m\ge1$ with
\begin{equation}\label{eq:criterion}
 c=\frac1m,\qquad A_k\mid mB_k\quad(k\ge0).
\end{equation}
When these conditions hold, the unique factor is
\begin{equation}\label{eq:digit-remainder}
 \nu_{A,B,m}=\Law\left(
 \sum_{k\ge0}\frac{J_k-(M_k-1)/2}{mB_k}\right),\qquad
 M_k=\frac{mB_k}{A_k},
\end{equation}
where the $J_k$ are independent and uniform on
$\{0,\ldots,M_k-1\}$. It is symmetric and has extreme support points
$\pm(R_A-R_B/m)$.
\end{theorem}
\begin{proof}
Assume a factor exists. Its compactness and the entire Fourier identity
follow from Lemma~\ref{lem:compact}. The source transform
$\Phi_B(cz)$ vanishes at $z=1/c$. Every target zero is a nonzero integer,
so $1/c=m$ is a positive integer.

Fix $k\ge0$ and evaluate zero orders at $z=mB_k$. In the source,
\[
 \ord_{z=mB_k}\Phi_B(z/m)=v_B(B_k)=k+1,
\]
because the divisors $B_j$ of $B_k$ are exactly those with $j\le k$.
The target must have order at least $k+1$. For a divisibility chain,
$\{j:A_j\mid mB_k\}$ is an initial segment. Having at least $k+1$
elements is therefore equivalent to $A_k\mid mB_k$. This proves all
conditions in \eqref{eq:criterion}.

Conversely suppose those conditions hold. The integers $M_k$ in
\eqref{eq:digit-remainder} are positive. Apply
Lemma~\ref{lem:uniform-split} to the target uniform of length $1/A_k$,
with small length $1/(mB_k)$ and digit count $M_k$. Construct all small
uniforms and all digits independently. Each target summand is then the
sum of one small uniform and one centered digit, and these target summands
are independent across $k$.

The digit series converges absolutely, since its $k$th term has absolute
value at most
\[
 \frac{M_k-1}{2mB_k}
 =\frac12\left(\frac1{A_k}-\frac1{mB_k}\right),
\]
and the sum of these bounds is finite. Summing the independent small
uniforms gives law $\D_{1/m}\mu_B$; summing the digits gives
\eqref{eq:digit-remainder}. This proves the factorization by an actual
probability measure, without an appeal to a merely formal quotient.

Uniqueness and symmetry follow from Lemma~\ref{lem:compact}. Its support
extrema must add to those of the source to give $\pm R_A$, so they are
$\pm(R_A-R_B/m)$. This conclusion concerns extreme points; the intervening
support need not be an interval.
\end{proof}

\begin{remark}[Why the indexing matters]
The argument would fail for a general unordered list of integer
frequencies: zero counts need not then form initial segments. The ladder
hypothesis is the mechanism that converts a zero-order inequality into a
coordinatewise divisibility condition. Positivity is supplied separately
by interval tiling. Neither step may be omitted in an attempted
extension to arbitrary products.
\end{remark}

\begin{corollary}[Equivalent zero-order criterion]\label{cor:zero-criterion}
For $m\ge1$, the factor in Theorem~\ref{thm:ladder} exists if and only if
\begin{equation}\label{eq:all-zero-criterion}
 v_A(mn)\ge v_B(n)\qquad(n\ge1).
\end{equation}
It is enough to test the generator frequencies $n=B_k$.
\end{corollary}
\begin{proof}
Necessity follows from the entire factorization. If all generator
inequalities hold, the proof of Theorem~\ref{thm:ladder} gives the
coordinatewise criterion. Conversely, if $B_j\mid n$ then
$A_j\mid mB_j\mid mn$, proving \eqref{eq:all-zero-criterion}.
\end{proof}

\subsection{The scale set is controlled by a single lcm obstruction}
\begin{definition}[Finite obstruction lcm]
Put $d_k=A_k/\gcd(A_k,B_k)$. If the integers $d_k$ have a common positive
integer multiple, let $L(A,B)$ be their least common multiple. Otherwise
write $L(A,B)=\infty$; the symbol $\infty$ here means absence of an ordinary
finite integer common multiple, not a permitted scale parameter.
\end{definition}

\begin{corollary}[All scales and a largest one]\label{cor:lcm}
If $L(A,B)<\infty$, all possible positive scales are exactly
\begin{equation}\label{eq:lcm-scales}
 \left\{\frac1{L(A,B)r}:r\in\N_{>0}\right\}.
\end{equation}
If $L(A,B)=\infty$, there is no possible positive scale. Equivalently,
finite existence requires a positive integer whose prime exponents dominate
\begin{equation}\label{eq:prime-deficits}
 \lambda_p=\sup_{k\ge0}\max\{0,v_p(A_k)-v_p(B_k)\}
\end{equation}
for every prime $p$. Thus every $\lambda_p$ must be finite and only finitely
many of them may be positive; then $L(A,B)=\prod_p p^{\lambda_p}$.
\end{corollary}
\begin{proof}
The elementary equivalence $A_k\mid mB_k$ if and only if
$A_k/\gcd(A_k,B_k)\mid m$ reduces Theorem~\ref{thm:ladder} to the definition
of the least common multiple. Prime factorization gives the last assertion.
Both finiteness requirements are essential: bounded deficits at each prime
do not help when infinitely many different primes are needed.
\end{proof}

A failed fixed-scale factorization has a finite certificate: a single index
$k$ and prime $p$ for which $v_p(A_k)>v_p(m)+v_p(B_k)$.
The real number $mB_k$ is then an explicit analytic obstruction, because the
source has too high a zero multiplicity there. In contrast, confirming the
criterion for arbitrarily many initial indices does not certify its truth
for all indices. Section~\ref{sec:computability} makes this distinction exact.

\begin{proposition}[Recovery and composition]\label{prop:composition}
The law $\mu_A$ determines its ladder, by
\begin{equation}\label{eq:recover-ladder}
 A_k=\min\{n\ge1:v_A(n)\ge k+1\}.
\end{equation}
If $A_k\mid mB_k$ and $B_k\mid nC_k$ for every $k$, then
\begin{equation}\label{eq:cocycle}
 \nu_{A,C,mn}=\nu_{A,B,m}*\D_{1/m}\nu_{B,C,n}.
\end{equation}
\end{proposition}
\begin{proof}
The integers with at least $k+1$ zero factors are exactly the positive
multiples of $A_k$, proving \eqref{eq:recover-ladder}. For composition,
substitute $\mu_B=\D_{1/n}\mu_C*\nu_{B,C,n}$ into the factorization of
$\mu_A$. The resulting remainder must be the unique factor for scale
$1/(mn)$, by Theorem~\ref{thm:ladder}.
\end{proof}

\subsection{Moment information carried by the factor}
Let $B_{2j}$ denote the usual even Bernoulli numbers, defined by
$z/(e^z-1)=\sum B_nz^n/n!$ near zero. For a compact law, its cumulants are
the derivatives at zero of the logarithm of its moment generating function.
\begin{proposition}[Exact cumulants]\label{prop:cumulants}
For every $j\ge1$,
\begin{align}
 \kappa_{2j}(\mu_A)&=\frac{B_{2j}}{2j}\sum_{k\ge0}A_k^{-2j},
 \label{eq:cumulants-law}\\
 \kappa_{2j}(\nu_{A,B,m})&=\frac{B_{2j}}{2j}
 \left(\sum_{k\ge0}A_k^{-2j}-m^{-2j}\sum_{k\ge0}B_k^{-2j}\right).
 \label{eq:cumulants-noise}
\end{align}
All odd cumulants vanish. In particular,
$\Var(X_A)=\frac1{12}\sum A_k^{-2}$.
\end{proposition}
\begin{proof}
For a centered unit uniform the moment generating function is
$\sinh(z/2)/(z/2)$. Differentiating its logarithm gives
$\frac12\coth(z/2)-1/z$, whose power series integrates to
\[
 \log\frac{\sinh(z/2)}{z/2}
 =\sum_{j\ge1}\frac{B_{2j}}{2j(2j)!}z^{2j}.
\]
Substitute $z/A_k$ and sum near zero. Normal convergence justifies the
coefficient calculation. Cumulants add under independent convolution and
scale homogeneously, yielding \eqref{eq:cumulants-noise}.
\end{proof}


\subsection{A quantitative obstruction, without a bound on the remainder}
For probability laws with finite first moments, write $W_1(\sigma,\tau)$
for the infimum of $\E|X-Y|$ over all couplings with these marginal laws.
A simple characteristic-function inequality already gives a useful
support-independent obstruction.

\begin{proposition}[A Fourier lower bound for approximate factors]
\label{prop:wasserstein}
For probability laws $\mu,\rho$ and every real $t\ne0$,
\begin{equation}\label{eq:fourier-wasserstein}
 W_1(\mu,\rho*\nu)\ge
 \frac{\bigl(|\widehat\mu(t)|-|\widehat\rho(t)|\bigr)_+}{2\pi|t|}
\end{equation}
for every probability remainder $\nu$ for which the left side is finite.
Suppose in particular that $A_k\nmid mB_k$, put $t_0=mB_k$ and
$r=v_A(t_0)<k+1$, and define
\begin{align}
 a_r&=\frac{|\Phi_A^{(r)}(t_0)|}{r!}>0,\notag\\
 D_r&=\frac{(2\pi R_A)^{r+1}+(2\pi R_B/m)^{r+1}}{(r+1)!}.
 \label{eq:quantitative-constants}
\end{align}
For every $0<h\le a_r/(2D_r)$ and every probability $\nu$,
\begin{equation}\label{eq:quantitative-obstruction}
 W_1(\mu_A,\D_{1/m}\mu_B*\nu)
 \ge\frac{a_r h^r}{4\pi(t_0+h)}>0,
\end{equation}
interpreting the bound trivially if the distance is infinite.
\end{proposition}
\begin{proof}
For a coupling, the inequality
$|e^{2\pi itx}-e^{2\pi ity}|\le2\pi|t|\,|x-y|$
gives
$|\widehat\mu(t)-\widehat{\rho*\nu}(t)|\le2\pi|t|W_1(\mu,\rho*\nu)$.
Also $|\widehat{\rho*\nu}(t)|\le|\widehat\rho(t)|$ because a probability
characteristic function has modulus at most one. The triangle inequality
proves \eqref{eq:fourier-wasserstein}.

At $t_0$ the target has order $r$, while the source has order at least
$r+1$. A compact law of radius $R$ has Fourier derivatives bounded on the
real axis by $(2\pi R)^j$ in order $j$. Taylor's formula with an integral
remainder therefore gives
\[
 |\Phi_A(t_0+h)|-|\Phi_B((t_0+h)/m)|
 \ge a_rh^r-D_rh^{r+1}\ge\frac12a_rh^r.
\]
Apply \eqref{eq:fourier-wasserstein} at $t=t_0+h$ to obtain the claimed
bound. No derivative, support, or moment estimate for $\nu$ was needed.
\end{proof}

The leading coefficient in this bound can itself be written explicitly:
\begin{equation}\label{eq:leading-coefficient}
 a_r=t_0^{-r}\prod_{j\ge r}|\sinc(t_0/A_j)|.
\end{equation}
Indeed, each of the $r$ vanishing factors has first derivative of absolute
value $1/t_0$ at $t_0$, and the remaining product is nonzero. Thus an
arithmetic failure yields not only an exact nonexistence witness but a
strict lower bound against \emph{every} independent probability remainder.
The bound may be very small for a high-index obstruction; no uniform
separation over all failed instances is implied.

\section{Geometric laws and cross-base rigidity}
\label{sec:geometric}

\subsection{All self-factoring contractions, for every real ratio}
\begin{theorem}[Geometric self-factor spectrum]\label{thm:self-spectrum}
For $0<q<1$,
\begin{equation}\label{eq:self-spectrum}
 \{c>0:\mu_q=\D_c\mu_q*\nu\text{ for some probability }\nu\}
 =\{q^j/m:j\in\N_0,\ m\in\N_{>0}\}.
\end{equation}
\end{theorem}
\begin{proof}
If the factor exists, its compactness forces the zero $1/c$ of the scaled
source to be a zero of the target. Lemma~\ref{lem:zeros} gives
$1/c=mq^{-j}$ with $m\ge1$ and $j\ge0$.

For the converse, first split the series at index $j$:
\[
 \mu_q=(\Unif_1*\Unif_q*\cdots*\Unif_{q^{j-1}})*\D_{q^j}\mu_q,
\]
where the first convolution is $\delta_0$ when $j=0$. Split every uniform
in the tail using the same integer digit count $m$, by
Lemma~\ref{lem:uniform-split}. The smaller uniforms sum to law
$\D_{q^j/m}\mu_q$, and the initial block and all the independent digit
terms form the probability remainder. Absolute convergence follows from
the geometric bound.
\end{proof}

This spectrum is countable and has no accumulation point other than zero:
for a fixed $\varepsilon>0$, $q^j/m\ge\varepsilon$ bounds both $j$ and $m$.
In particular, these laws do not admit such a self-factor decomposition at
every contraction $0<c<1$. For integer reciprocal $q=1/a$, the spectrum
simplifies to all reciprocal positive integers, consistent with
Theorem~\ref{thm:ladder} for $A=B$.

\subsection{An integer-base target forces an integer-base source}
\begin{theorem}[Cross-base classification]\label{thm:crossbase}
Let $a\ge2$ be an integer, $0<q<1$, and $c>0$. Then
\begin{equation}\label{eq:crossbase}
 \Mu a=\D_c\mu_q*\nu
\end{equation}
for a probability $\nu$ if and only if
\begin{equation}\label{eq:crossbase-criterion}
 q=\frac1b\text{ for an integer }b\ge2,\qquad a\mid b,\qquad
 c=\frac1m\text{ for an integer }m\ge1.
\end{equation}
When this holds, the digit count in \eqref{eq:digit-remainder} is
$M_k=m(b/a)^k$.
\end{theorem}
\begin{proof}
The zero at $1/c$ first gives $c=1/m$. Every zero $mq^{-k}$ of the source
must be an integer zero of the target. Write $b=1/q$. Since $mb$ is an
integer, $b$ is rational. In lowest terms let $b=u/v$ with positive coprime
integers $u,v$. The integrality of $mu^k/v^k$ for every $k$ forces
$v^k\mid m$ for every $k$, which is impossible unless $v=1$.
Thus $b$ is an integer at least two.

Theorem~\ref{thm:ladder} now says $a^k\mid mb^k$ for every $k$.
For each prime this is
$k v_p(a)\le v_p(m)+k v_p(b)$. Letting $k$ increase forces
$v_p(a)\le v_p(b)$, hence $a\mid b$. Conversely that divisibility makes
$a^k\mid mb^k$ automatic, and the explicit construction applies.
\end{proof}

\begin{corollary}[The Rvachev target]
A geometric-uniform law can be a positively scaled convolution factor of
$\Mu2$ exactly when its ratio is $1/b$ with $b$ even. The permitted scales
are precisely $1/m$, $m\ge1$. In particular, bases $6$ and $10$ are allowed,
not only powers of $2$; base $3$ and all irrational ratios are excluded.
\end{corollary}

\begin{example}[Support and variance do not decide factorability]
Consider the target $\Mu3$, the source $\Mu2$, and scale $c=1/2$.
The target has support radius $3/4$ and variance $3/32$; the scaled source
has radius $1/2$ and variance $1/36$. These necessary size inequalities are
satisfied. Nevertheless, at $t=4$ the source transform $\Phi_{(2^k)}(t/2)$
has a zero of order $2$, whereas $\Phi_{(3^k)}(t)$ has order $1$.
No probability remainder exists. This obstruction survives arbitrarily
strong shrinking by changing the frequency at which the contradiction is
read; the failure $3\nmid2$ rules out every positive scale.
\end{example}

\begin{example}[A non-dyadic factor of the dyadic law]
For target base $2$, source base $6$, and $m=1$, the complement is
\[
 \Law\left(\sum_{k\ge0}
 \frac{J_k-(3^k-1)/2}{6^k}\right),\qquad
 J_k\text{ uniform on }\{0,\ldots,3^k-1\}.
\]
The $k=0$ digit is deterministic. The support radius of this remainder is
$1-3/5=2/5$. The theorem proves its existence and uniqueness, but does not
by itself determine its absolute continuity or exact smoothness.
\end{example}

\section{A compact smooth realization of inclusion modulo finite sets}
\label{sec:boolean}

\subsection{Encoding one bit by one prime}
Let $p_1,p_2,\ldots$ list the odd primes in increasing order.
For $S\subseteq\N_{>0}$ define
\begin{equation}\label{eq:boolean-ladder}
 A_{S,0}=1,\qquad
 A_{S,k}=2^k\prod_{\substack{1\le j\le k\\j\in S}}p_j\quad(k\ge1),
 \qquad \mu_S=\mu_{A_S}.
\end{equation}
Each ratio is either $2$ or $2p_{k+1}$, so these are divisibility ladders.
Every $\mu_S$ consequently has a smooth density $f_S$ supported in $[-1,1]$.
In fact, the whole family consists of unscaled convolution divisors of the
single classical Rvachev law $\Mu2$.

Write $S\subseteq^*T$ when $S\setminus T$ is finite, and $S\mathrel{E_0}T$
when $S\triangle T$ is finite. The symbol $E_0$ denotes finite-change
 equivalence on binary sequences; its classical role in Borel
classification is discussed in \cite{hkl}. We will prove the obstruction
needed here directly.

\begin{theorem}[Exact arithmetic encoding]\label{thm:boolean}
For $S,T\subseteq\N_{>0}$ and $m\ge1$,
\begin{equation}\label{eq:boolean-factor}
 \mu_S=\D_{1/m}\mu_T*\nu
 \quad\Longleftrightarrow\quad
 S\setminus T\text{ is finite and }
 \prod_{j\in S\setminus T}p_j\mid m.
\end{equation}
In particular, an unscaled factorization exists exactly when $S\subseteq T$;
a factorization at some positive scale exists exactly when $S\subseteq^*T$.
Mutual factorability at possibly different positive scales is equivalent
to $S\mathrel{E_0}T$. The unscaled laws themselves are equal only when $S=T$.
Every $\mu_S$ is an unscaled convolution factor of $\Mu2$.
\end{theorem}
\begin{proof}
For each $k$ cancellation of the common factors gives
\[
 \frac{A_{S,k}}{\gcd(A_{S,k},A_{T,k})}
 =\prod_{\substack{j\le k\\j\in S\setminus T}}p_j.
\]
The least common multiple of these initial products is finite exactly when
$S\setminus T$ is finite, and is then the product in
\eqref{eq:boolean-factor}. Corollary~\ref{cor:lcm} proves the equivalence
and all the assertions about scales. Applying it in both directions gives
finite symmetric difference. Finally equality of laws recovers equality of
ladders by Proposition~\ref{prop:composition}; the ratio at index $j$
recovers membership of $j$ in $S$. Taking $S=\varnothing$ in the
unscaled criterion gives $\mu_{\varnothing}=\Mu2=\mu_T*\nu$ for
every $T$, proving the final assertion.
\end{proof}

\begin{remark}[Orientation of the order]
The displayed equations fix the order convention: the \emph{target}
$\mu_S$ contains a scaled copy of the \emph{source} $\mu_T$ when
$S\subseteq^*T$. Larger sets introduce more prime denominators and hence
smaller summands. This avoids the ambiguity of calling convolution
containment an order without specifying its direction.
\end{remark}

\subsection{The encoding is continuous in all derivatives}
Equip $2^{\N_{>0}}$ with its product topology, and equip
$C^\infty_{[-1,1]}(\R)$ with the seminorms
$\|f^{(r)}\|_\infty$, $r\ge0$, restricting to functions supported in
$[-1,1]$.

\begin{proposition}[A compact family of smooth densities]\label{prop:continuous}
The map $S\mapsto f_S$ is continuous and injective, and hence a
homeomorphism onto a compact subset $\mathcal K$ of
$C^\infty_{[-1,1]}(\R)$. More explicitly, put
\[
 \overline A_j=2^j\prod_{i=1}^jp_i,\qquad
 C_r=2^{r+1}\prod_{j=0}^{r+1}\overline A_j.
\]
If $S,T$ agree on the first $N$ positions and $N\ge r+2$, then
\begin{equation}\label{eq:continuity-bound}
 \|f_S^{(r)}-f_T^{(r)}\|_\infty\le C_r2^{-N}.
\end{equation}
\end{proposition}
\begin{proof}
Their prefix densities through index $N$ agree. Couple the remaining tails
using the same uniforms. Each tail is bounded by $1/(2A_{S,N})$, and
$A_{S,N}=A_{T,N}\ge2^N$. Therefore the difference of the two tails has
absolute value at most $1/A_{S,N}\le2^{-N}$.
The common prefix has $r$th derivative Lipschitz with constant
$L_r(A_S)\le C_r$, by the proof of
Proposition~\ref{prop:approximation}. Taking expectations of the shifted
prefix proves \eqref{eq:continuity-bound}. This is continuity in every
seminorm. Injectivity was proved in Theorem~\ref{thm:boolean}; a continuous
injection from compact Cantor space into a Hausdorff space is a
homeomorphism onto its image.
\end{proof}

The common compact support and the continuity in every derivative are
important. The classification barrier below is not produced by a sequence
of increasingly singular target densities, by unbounded supports, or by
a discontinuous coding map. It is present inside one compact family of
ordinary smooth probability densities.

\subsection{No Borel real-valued complete invariant}
Define $f\sim g$ on $\mathcal K$ when each of the corresponding laws is a
convolution of a positive dilation of the other with a probability measure.
A real-valued \emph{complete invariant} would be a map $I$ such that
$I(f)=I(g)$ holds exactly when $f\sim g$. The requirement that $I$ be
Borel measurable excludes arbitrary nonmeasurable choices of representatives.

\begin{theorem}[Nonsmooth classification on a compact smooth family]
\label{thm:nonsmooth}
There is no Borel measurable real-valued complete invariant for $\sim$ on
$\mathcal K$. There is also no complete invariant consisting of countably
many Borel real-valued coordinates. In fact, every continuous real-valued
invariant for $\sim$ is constant on $\mathcal K$.
\end{theorem}
\begin{proof}
Suppose first that $I$ were a Borel complete invariant, and pull it back
along the continuous map of Proposition~\ref{prop:continuous}:
$J(S)=I(f_S)$. By Theorem~\ref{thm:boolean}, $J$ is unchanged under every
finite modification of $S$.

Give the bits of $S$ the independent fair-coin law. For each rational $t$,
the Borel event $\{J\le t\}$ is unchanged by modifying the first $N$ bits,
for every $N$. It is therefore measurable with respect to the bits after
$N$: fix the first $N$ bits to zero to see that its membership depends only
on the remaining coordinates. Hence it is a tail event.

A tail event has probability zero or one. One elementary proof is that
such an event is independent of each finite-coordinate cylinder; the
monotone-class argument extends independence to the sigma-algebra generated
by all cylinders, which contains the event itself. Its probability $p$
then obeys $p=p^2$.
It follows that the distribution function of the real random variable
$J$ takes only the values zero and one at rational arguments. Consequently
$J$ is almost surely equal to a constant: take the cut between rational
arguments of probability zero and those of probability one. The cut is
finite because $J$ is real valued, and continuity of probability from
above and below gives the claim.

Completeness of $I$ would place this conull level set inside a single
$E_0$-class. But each such class is countable, since there are only
countably many finite changes of a fixed sequence; each individual sequence
has fair-coin measure zero. This is a contradiction.

For countably many Borel real coordinates, apply the same argument to each
coordinate and intersect their conull sets of constancy. The complete
coordinate vector is again constant on a conull set, giving the same
contradiction. Finally each $E_0$-class is dense in Cantor space, since a
finite change can prescribe any finite initial block. A continuous invariant
constant on one such dense class is constant everywhere.
\end{proof}

\begin{remark}[What the obstruction does not say]
The full Fourier transform remains a complete invariant for \emph{equality}
of probability laws. The theorem instead concerns the coarser relation
of mutual scaled factorability. It does not say that this relation is
non-Borel: on the constructed family it is the familiar Borel relation
$E_0$. Nor does it prohibit a nonmeasurable choice-based assignment of real
labels. The precise negative conclusion is the absence of measurable
complete labels by equality of real numbers or countable real sequences.
\end{remark}

\section{Effective classification: exact arithmetical complexity}
\label{sec:computability}

\subsection{Input model and the meaning of completeness}
An effective assertion requires a presentation model. Use finite program
codes for primitive-recursive functions specifying the positive ladder
ratios, with validity guaranteed by construction. For example, from any
primitive-recursive $h:\N_0\to\N_0$ take
\[
 A_0=1,\qquad A_{k+1}=A_k(2+h(k)).
\]
This represents valid ladders without a separate undecidable totality or
validity test. All the hardness constructions below have ratios of this
form. The conclusions also hold as promise-problem statements for any
larger presentation system containing these examples and providing their
integer entries effectively.

A predicate is $\Pi^0_1$ when it has the form $\forall k\,R(e,k)$ with
a decidable relation $R$; it is $\Sigma^0_2$ when it has the form
$\exists m\,\forall k\,R(e,m,k)$. Completeness means that every predicate
at the indicated level can be transformed effectively into the problem
under consideration, preserving yes and no answers. For the hardness
proofs we use the usual primitive-recursive matrix normal forms, so the
constructed ladder entries remain primitive recursive.

\begin{theorem}[Effective complexity of exact factors]\label{thm:complexity}
For primitive-recursively presented divisibility ladders $A,B$:
\begin{enumerate}[label=(\roman*)]
\item For each fixed integer $m\ge1$, existence of a factor in
$\mu_A=\D_{1/m}\mu_B*\nu$ is $\Pi^0_1$-complete.
\item Existence of a factor at \emph{some} positive real scale $c$ is
$\Sigma^0_2$-complete.
\end{enumerate}
The target and source densities in the hardness examples are uniformly
computable together with all derivatives, with the quantitative error
bounds of Proposition~\ref{prop:approximation}.
\end{theorem}
\begin{proof}
For (i), the upper bound follows from the criterion
$\forall k\;(A_k\mid mB_k)$, which is a universal quantification of a
decidable integer test.

For hardness, start with an arbitrary $\Pi^0_1$ predicate
$\forall t\,R(e,t)$ in primitive-recursive normal form. Choose an odd prime
$p$ that does not divide the fixed integer $m$. Set $\varepsilon_0=0$ and,
for $k\ge1$, set $\varepsilon_k=1$ if some $t<k$ violates $R(e,t)$,
and $\varepsilon_k=0$ otherwise. This is a nondecreasing binary sequence
computed by a bounded search. Define
\[
 B_k=2^k,\qquad A_k=2^kp^{\varepsilon_k}.
\]
Both are valid primitive-recursive ladders. If there is no violation, they
are equal and the factor exists. If there is a violation, all sufficiently
late $A_k$ contain the extra prime $p$, which cannot divide $mB_k$.
The factor therefore exists exactly when the original $\Pi^0_1$
predicate holds.

For (ii), Theorem~\ref{thm:ladder} eliminates quantification over arbitrary
real scales and gives the $\Sigma^0_2$ upper bound
\[
 \exists m\ge1\ \forall k\ge0\;(A_k\mid mB_k).
\]
To prove hardness, let $W_e$ be a computably enumerable set of positive
integers, with a fixed bounded-stage enumeration. Define
\begin{equation}\label{eq:ce-ladders}
 B_k=2^k,\qquad
 A_k=2^k\prod_{\substack{1\le j\le k\\
                    j\text{ enumerated in }W_e\text{ by stage }k}}p_j.
\end{equation}
Membership in this finite stage condition is decided by bounded simulation;
products and the enumeration of the first $k$ odd primes are primitive
recursive. The set of prime factors included is nondecreasing in $k$, so
these are valid ladders. Every $j\in W_e$ eventually contributes $p_j$.
Corollary~\ref{cor:lcm} shows that some scale exists if and only if $W_e$
is finite.

Here is a direct verification of the needed hardness of finiteness, rather
than an unsupported appeal to it. Given
$\exists n\,\forall t\,R(e,n,t)$ with primitive-recursive $R$, let
$\ell_s$ be the least $n\le s$ for which $R(e,n,t)$ holds for every
$t\le s$, or $s+1$ if no such $n$ exists. Enumerate a fresh integer whenever
$\ell_s$ changes. If there is a true witness, the least one eventually
survives while all smaller candidates have been refuted, so $\ell_s$
stabilizes and the enumerated set is finite. If there is no witness, every
bounded collection of candidates is eventually refuted, so
$\ell_s\to\infty$ and changes infinitely often. Thus the enumerated set
is finite exactly when the given $\Sigma^0_2$ predicate holds. Composing
this construction with \eqref{eq:ce-ladders} proves hardness.

Finally, every entry of the constructed ladders is uniformly computable
from the input code. The rational spline and explicit derivative error
bound in Proposition~\ref{prop:approximation} apply uniformly to them.
The complexity is therefore a property of exact factorization, not a
failure to compute the smooth input functions.
\end{proof}

\subsection{Consequences for certificate-based computation}
For a fixed reciprocal-integer scale, a search for a failed divisibility
condition semidecides nonexistence: the first failed index is an exact
certificate. A general terminating algorithm for both existence and
nonexistence would contradict Theorem~\ref{thm:complexity}. Likewise,
searching for successively better numerical agreement between transforms
cannot solve the unrestricted exact problem.

This does not obstruct useful algorithms on structured subclasses. For
constant integer bases, Theorem~\ref{thm:crossbase} reduces everything to
finite prime factorization. More generally, a symbolic description that
proves eventual periodicity or bounded prime deficits can certify the
infinite condition. The theorem identifies what additional structural
information an implementation must exploit; it does not say that every
individual instance is difficult.

There is also a separation between computing a factor and recognizing
that it exists. Once the arithmetic hypotheses are provided with a proof,
\eqref{eq:digit-remainder} computes the factor weakly by its finite digit
prefixes with an explicit tail bound. Such a computation is perfectly
compatible with undecidability of the unpromised existence problem.


\subsection{A terminating algorithm for eventually periodic ratios}
\begin{proposition}[Periodic-ratio decision theorem]\label{prop:periodic}
Suppose the integer ratios $A_{k+1}/A_k$ and $B_{k+1}/B_k$ are periodic
for $k\ge K$ with a common period $P\ge1$. Put
\[
 Q_A=A_{K+P}/A_K,\qquad Q_B=B_{K+P}/B_K.
\]
Some positive-scale factor exists if and only if $Q_A\mid Q_B$.
When this holds,
\begin{equation}\label{eq:periodic-lcm}
 L(A,B)=\lcm_{0\le k<K+P}\frac{A_k}{\gcd(A_k,B_k)}.
\end{equation}
Thus the existence problem, all scales, and the largest scale are decided
by a finite calculation from a certified eventually periodic presentation.
\end{proposition}
\begin{proof}
For $0\le r<P$ and $j\ge0$, periodicity gives
$A_{K+jP+r}=A_{K+r}Q_A^j$ and
$B_{K+jP+r}=B_{K+r}Q_B^j$.
At any prime $p$ the exponent difference is therefore
\[
 v_p(A_{K+r})-v_p(B_{K+r})
   +j\bigl(v_p(Q_A)-v_p(Q_B)\bigr).
\]
A positive slope makes the deficits unbounded and excludes every finite
$m$. When all slopes are nonpositive, their maxima are attained at $j=0$;
together with the initial indices these are exactly the finite tests in
\eqref{eq:periodic-lcm}. Corollary~\ref{cor:lcm} completes the proof.
\end{proof}

For example, alternating ratios $(2,6)$ and $(3,4)$ have the same period
product $12$. Their laws are mutually scaled-factorable, with minimal
scale denominators $2$ and $3$ in the respective directions. On the other
hand, period products $12$ and $18$ do not admit the factor in the first
direction: ordinary numerical size $12<18$ is insufficient because
$12\nmid18$. This finite algorithm and the unrestricted impossibility
result apply to different input classes and are entirely consistent.

\section{Same-base remainders: singularity and finite smoothing}
\label{sec:regularity}

\subsection{Digit representation and extraction of full base factors}
Fix integers $a\ge2$ and $m\ge1$. Let $\eta_{a,m}$ be the unique measure
such that
\begin{equation}\label{eq:same-base}
 \Mu a=\D_{1/m}\Mu a*\eta_{a,m}.
\end{equation}
Theorem~\ref{thm:ladder} gives
\begin{equation}\label{eq:same-digits}
 \eta_{a,m}=\Law\left(\frac1m\sum_{k\ge0}a^{-k}
       \left(J_k-\frac{m-1}{2}\right)\right),
 \qquad J_k\text{ uniform on }\{0,\ldots,m-1\}.
\end{equation}
The support radius is $(1-1/m)a/(2(a-1))$. For $m=1$ this is $\delta_0$.

Write
\begin{equation}\label{eq:base-valuation}
 m=a^d r,\qquad d=\max\{j\ge0:a^j\mid m\},\qquad a\nmid r.
\end{equation}
When $r=1$, the last condition holds automatically because $a\ge2$.
For composite $a$, $d$ is divisibility by powers of the \emph{whole} base;
it is not the valuation at a single arbitrarily chosen prime.

\begin{proposition}[Uniform smoothing decomposition]\label{prop:smoothing}
With \eqref{eq:base-valuation},
\begin{equation}\label{eq:smoothing}
 \eta_{a,m}=
 \left(\mathop{*}_{j=0}^{d-1}\Unif_{a^{-j}}\right)
 *\D_{a^{-d}}\eta_{a,r},
\end{equation}
where an empty convolution is $\delta_0$.
\end{proposition}
\begin{proof}
The self-factor formula for the scale $a^{-d}$ is obtained by splitting
the geometric series into its first $d$ terms and its tail. Thus
$\eta_{a,a^d}=\mathop{*}_{j<d}\Unif_{a^{-j}}$.
Apply the composition identity \eqref{eq:cocycle} with scales $a^{-d}$
and $1/r$. This gives \eqref{eq:smoothing} without division at Fourier zeros.
Equivalently, the quotient of transforms telescopes through the first $d$
sinc factors.
\end{proof}

\subsection{The primitive remainder is singular continuous}
Define the centered digit mask
\begin{equation}\label{eq:mask}
 h_r(u)=\frac1r\sum_{j=0}^{r-1}e^{2\pi iu(j-(r-1)/2)}.
\end{equation}
The absolute value $|h_r|$ is one-periodic. The mask itself equals
$h_r(u)=\sin(\pi ru)/(r\sin(\pi u))$ whenever the denominator is nonzero,
with the removable values filled in by \eqref{eq:mask}. Thus
\begin{equation}\label{eq:eta-transform}
 \widehat\eta_{a,r}(t)=\prod_{k\ge0}h_r\left(\frac{t}{ra^k}\right).
\end{equation}

\begin{lemma}[Exact Fourier resonance]\label{lem:resonance}
If $r\ge2$ and $a\nmid r$, the number
\begin{equation}\label{eq:resonance-constant}
 C_{a,r}:=\prod_{\ell\ge1}
 \left|\frac{\sin(\pi r/a^\ell)}{r\sin(\pi/a^\ell)}\right|
\end{equation}
is positive, and for every $N\ge0$,
\begin{equation}\label{eq:resonance}
 |\widehat\eta_{a,r}(ra^N)|=C_{a,r}.
\end{equation}
\end{lemma}
\begin{proof}
At $t=ra^N$, the mask arguments for $k\le N$ are integers and their
absolute values are one. Reindexing the remaining factors gives
\eqref{eq:resonance-constant}. Its denominators are nonzero. A numerator
could vanish only if $a^\ell\mid r$, which would imply $a\mid r$.
Every factor is therefore nonzero. As $\ell\to\infty$ the factor is
$1+O(a^{-2\ell})$, so the product of their absolute values converges to
a nonzero number.
\end{proof}

\begin{theorem}[Primitive digit remainder]\label{thm:primitive}
If $r\ge2$ and $a\nmid r$, the probability measure $\eta_{a,r}$ is
atomless and singular with respect to Lebesgue measure. If $r<a$, its
self-similar support has disjoint first-level pieces. If $r\ge a$, its
support is its entire closed support interval, despite the singularity.
\end{theorem}
\begin{proof}
The random series satisfies the self-similar distributional equation
\begin{equation}\label{eq:ifs}
 Y\ \stackrel{d}{=}\ \frac{Y'}a+\frac{J-(r-1)/2}{r},
\end{equation}
with $Y'$ an independent copy and $J$ a uniform digit. This equation has a
unique probability solution. Indeed, iterate its characteristic-function
equation $n$ times; the remaining characteristic factor is evaluated at
$t/a^n$ and tends to one. The other factors converge to
\eqref{eq:eta-transform}, determining the law uniquely. This is a direct
special-case proof of the contracting self-similar measure construction
associated with \cite{hutchinson}.

The affine maps in \eqref{eq:ifs} preserve absolute continuity and also
preserve singularity of measures with respect to Lebesgue measure. The
operator that averages their pushforwards therefore preserves both parts
of the Lebesgue decomposition separately. If the stationary measure had
both parts with positive mass, each part after normalization would be a
stationary probability, contradicting uniqueness. The measure is consequently
of pure type. Its Fourier transform does not tend to zero by
Lemma~\ref{lem:resonance}, whereas the Fourier transform of every integrable
density tends to zero. The latter elementary Riemann--Lebesgue assertion
follows by approximating the density in $L^1$ by step functions on bounded
intervals, whose Fourier integrals tend to zero explicitly. The measure
cannot be absolutely continuous, and hence is purely singular.

To exclude atoms, suppose the largest atom mass were $p>0$. Such a maximum
is attained: only finitely many atoms can have mass at least half the
supremum. Let $E$ be the finite nonempty set of atoms of mass $p$.
For $x\in E$, the average in \eqref{eq:ifs} can equal $p$ only if every
inverse image under all $r$ affine maps also has mass $p$. Each inverse
map is injective and sends $E$ into $E$, hence permutes it. An affine map
with expansion factor $a>1$ cannot permute a finite set of positive diameter,
so $E$ is a singleton. Distinct digit maps have distinct fixed points and
cannot all preserve that singleton. This contradiction proves atomlessness.

Finally the convex hull of the support has length
$H=(r-1)a/(r(a-1))$. Its first-level images have length $H/a$ and adjacent
left endpoints separated by $1/r$. They are disjoint when $r<a$ and
cover the full hull when $r\ge a$. In the latter case the hull is an
invariant compact set, so it equals the self-similar attractor. Uniqueness
of that compact attractor follows by iterating the maps: their Hausdorff
contraction ratio is $1/a$. Alternatively, the possible digit prefixes
plus a tail of diameter $Ha^{-n}$ converge to the same set.
\end{proof}

\subsection{Sharp differentiability of every same-base remainder}
The space $W^{d,1}(\R)$ consists of integrable functions whose distributional
derivatives through order $d$ are integrable functions. A function is of
bounded variation when its distributional first derivative is a finite
signed measure. These distinctions matter for a continuous function whose
derivative is singular, as with a Cantor distribution function.

\begin{theorem}[Complete integer regularity classification]
\label{thm:regularity}
Write $m=a^dr$ as in \eqref{eq:base-valuation}.
\begin{enumerate}[label=(\roman*)]
\item If $m=1$, then $\eta_{a,m}=\delta_0$.
\item If $d=0$ and $r\ge2$, then $\eta_{a,m}$ is singular continuous.
\item If $d\ge1$ and $r=1$, then $\eta_{a,m}$ is the convolution of the
$d$ uniforms in \eqref{eq:smoothing}. For $d=1$ its density has no continuous
representative. For $d\ge2$ its density is $C^{d-2}$ but not $C^{d-1}$.
In every case it has bounded weak derivatives through order $d-1$.
\item If $d\ge1$ and $r\ge2$, then $\eta_{a,m}$ has a compactly supported
density $f\in C^{d-1}(\R)$. The derivative $f^{(d-1)}$ is continuous and
of bounded variation, but is not absolutely continuous. Consequently
$f\notin W^{d,1}(\R)$ and $f\notin C^d(\R)$.
\end{enumerate}
\end{theorem}
\begin{proof}
The first two assertions follow immediately from the digit representation
and Theorem~\ref{thm:primitive}. For (iii), repeated convolution of interval
indicators gives a degree-$(d-1)$ spline, with continuous derivatives
through order $d-2$ when $d\ge2$. Differentiate $d-1$ different uniform
factors distributionally and leave the last uniform undifferentiated;
this also shows that the $(d-1)$st weak derivative is bounded.
Near the left endpoint $L$ of the support, before any other spline knot,
the density is a positive constant times $(x-L)_+^{d-1}$.
Its $(d-1)$st derivative has a nonzero jump, proving the sharp failure
of $C^{d-1}$. For $d=1$ this is exactly the jump of a uniform density.

For (iv), put $\lambda=\D_{a^{-d}}\eta_{a,r}$, which is atomless singular.
After differentiating $d-1$ of the uniform factors, $f^{(d-1)}$ is a finite
linear combination of translates of the convolution of one uniform with
$\lambda$. Such a convolution is continuous: for a uniform of length
$\ell$ its density is
\[
 x\longmapsto \frac{\lambda([x-\ell/2,x+\ell/2])}{\ell},
\]
with endpoint conventions immaterial because $\lambda$ is atomless.
Its continuity follows from continuity of the distribution function of
$\lambda$. This proves $f\in C^{d-1}$.

Differentiating all $d$ uniform factors instead gives
\begin{equation}\label{eq:singular-derivative}
 D^df=\left(\mathop{*}_{j=0}^{d-1}
 a^j(\delta_{-a^{-j}/2}-\delta_{a^{-j}/2})\right)*\lambda,
\end{equation}
a finite signed sum of translates of a singular measure. It is a singular
finite signed measure. It is nonzero, since $D^df=0$ would make $f$ a
polynomial of degree at most $d-1$ as a distribution, and a compactly
supported polynomial must be zero, contradicting mass one. Thus
$f^{(d-1)}$ is of bounded variation but cannot be absolutely continuous:
its derivative is a nonzero singular measure. The asserted Sobolev and
classical nonmembership follow.
\end{proof}

\begin{center}
\begin{tabular}{@{}cc>{\raggedright\arraybackslash}p{9.0cm}@{}}
\toprule
Base $a$ & $m$ & Remainder type and sharp integer regularity\\
\midrule
2 & 2 & One uniform; discontinuous density.\\
2 & 3 & Singular continuous with interval support.\\
2 & 4 & Two-uniform spline; continuous, not $C^1$.\\
2 & 6 & Continuous density with a nonzero singular derivative.\\
2 & 8 & Three-uniform spline; $C^1$, not $C^2$.\\
2 & 12 & $C^1$ density whose second derivative is singular.\\
3 & 2 & Separated two-digit Cantor-type measure.\\
3 & 6 & Continuous density of exact H\"older exponent $\log 2/\log 3$.\\
\bottomrule
\end{tabular}
\end{center}
The last entry uses the sharper result proved in
Section~\ref{sec:holder}. The table concerns the remainder, not the target:
all geometric target laws in it have $C^\infty$ densities.

\section{Sharp Fourier order and separated-digit H\"older regularity}
\label{sec:fourier}

\subsection{The exact power envelope after smoothing}
The density regularity in Theorem~\ref{thm:regularity} has a matching
Fourier statement. The nondecaying primitive digit transform loses exactly
one power for each extracted uniform factor, not an unspecified amount.

\begin{lemma}[Offset resonance]\label{lem:offset}
Let $r\ge2$, $a\nmid r$, and $\delta\in\R$. Then
\begin{equation}\label{eq:offset}
 \lim_{N\to\infty}|\widehat\eta_{a,r}(ra^N+\delta)|
 =C_{a,r}|\widehat\eta_{a,r}(\delta)|.
\end{equation}
\end{lemma}
\begin{proof}
Split the mask product \eqref{eq:eta-transform} at $k=N$. For $k\le N$,
periodicity of the absolute value of $h_r$ gives
\[
 \left|h_r\left(a^{N-k}+\frac{\delta}{ra^k}\right)\right|
 =\left|h_r\left(\frac{\delta}{ra^k}\right)\right|.
\]
Their product converges to $|\widehat\eta_{a,r}(\delta)|$.
For $k=N+\ell$ with $\ell\ge1$, the factors are
\[
 \left|h_r\left(a^{-\ell}+\frac{\delta}{ra^{N+\ell}}\right)\right|.
\]
For each fixed $\ell$ they tend to $|h_r(a^{-\ell})|$. The convergence
of the infinite products is justified uniformly in $N$ after discarding
finitely many initial values of $N$: for large $\ell$, all arguments are
bounded by a constant times $a^{-\ell}$ and
$h_r(u)=1+O(u^2)$ near zero. The summable bound $O(a^{-2\ell})$ controls
the product tails uniformly. The limiting tail product is $C_{a,r}$.
\end{proof}

\begin{theorem}[Sharp power Fourier decay]\label{thm:fourier}
Let $m=a^dr$ as in \eqref{eq:base-valuation}, with $d\ge1$. Then
\begin{equation}\label{eq:fourier-order}
 |\widehat\eta_{a,m}(t)|=O(|t|^{-d})\qquad(|t|\to\infty),
 \qquad
 \limsup_{t\to\infty}t^d|\widehat\eta_{a,m}(t)|>0.
\end{equation}
More precisely, let $t_N=ma^N+1/2$. With the conventions
$C_{a,1}=1$ and $\widehat\eta_{a,1}=1$,
\begin{align}\label{eq:fourier-limit}
 \lim_{N\to\infty}t_N^d|\widehat\eta_{a,m}(t_N)|
 &=\frac{a^{d(d-1)/2}}{\pi^d}
   \left(\prod_{j=0}^{d-1}\sin\frac{\pi}{2a^j}\right)
   C_{a,r}\left|\widehat\eta_{a,r}\left(\frac1{2a^d}\right)\right|
 >0.
\end{align}
Thus the power $d$ cannot be replaced by any larger power in a global
upper bound. When $d=0$ and $r\ge2$, there is no decay to zero at all.
\end{theorem}
\begin{proof}
Taking Fourier transforms in \eqref{eq:smoothing} gives
\[
 \widehat\eta_{a,m}(t)
 =\left(\prod_{j=0}^{d-1}\sinc(t/a^j)\right)
   \widehat\eta_{a,r}(t/a^d).
\]
The last factor has modulus at most one, and the finite product is
$O(|t|^{-d})$.
For the special sequence $t_N$, every $ma^N/a^j$ with $j<d$ is an integer.
Consequently
\[
 t_N^d\prod_{j=0}^{d-1}|\sinc(t_N/a^j)|
 =\frac{a^{d(d-1)/2}}{\pi^d}
   \prod_{j=0}^{d-1}\sin\frac{\pi}{2a^j}.
\]
The remaining argument is
$t_N/a^d=ra^N+1/(2a^d)$. Apply Lemma~\ref{lem:offset}, or use the constant
transform if $r=1$, to obtain \eqref{eq:fourier-limit}.
Every finite sine factor is positive. The residual transform at
$\delta=1/(2a^d)$ is nonzero: in \eqref{eq:eta-transform} its numerator
arguments are $\delta/a^k\in(0,1)$, no factor vanishes, and the tail product
is nonzero. Finally the $d=0$ assertion is Lemma~\ref{lem:resonance}.
\end{proof}

A positive limsup is the appropriate sharpness statement, not a pointwise
lower bound: these transforms have many zeros. In particular, the theorem
does not assert $|\widehat\eta(t)|\asymp|t|^{-d}$ at every large frequency.

\subsection{An exact fractional exponent in the separated case}
\label{sec:holder}
When $2\le r<a$, let $\alpha=\log r/\log a\in(0,1)$. For a function $g$,
H\"older continuity of exponent $\alpha$ means
$|g(x)-g(y)|\le C|x-y|^\alpha$ for all real $x,y$ and a fixed constant $C$.

\begin{theorem}[Sharp separated-digit H\"older exponent]\label{thm:holder}
Assume $2\le r<a$.
The distribution function of $\eta_{a,r}$ is globally H\"older continuous
of exponent $\alpha=\log r/\log a$, and is not H\"older continuous with
any exponent $\beta$ satisfying $\alpha<\beta\le1$.
For $d\ge1$ and $m=a^dr$, the density in Theorem~\ref{thm:regularity}
satisfies
\begin{equation}\label{eq:holder}
 f\in C^{d-1,\alpha}(\R),\qquad
 f\notin C^{d-1,\beta}(\R)\quad(\alpha<\beta\le1).
\end{equation}
Here $C^{d-1,\alpha}$ means that the $(d-1)$st derivative is globally
H\"older of exponent $\alpha$.
\end{theorem}
\begin{proof}
Let $H$ be the convex-hull length of the primitive support. At level $n$,
there are $r^n$ cylinder sets, each of mass $r^{-n}$, each contained in an
interval of length $Ha^{-n}$. Strong separation gives a constant $g>0$
such that distinct level-$n$ hull intervals are separated by at least
$ga^{-(n-1)}$. Indeed, two cylinders first differing at level $j\le n$
lie in different first-level pieces of a common level-$(j-1)$ cylinder;
scaling the first-level gap gives at least this bound.

An interval of length at most $Ha^{-n}$ can therefore meet only a bounded
number of level-$n$ cylinders, independently of $n$. Choose $n$ so that
$Ha^{-(n+1)}<|I|\le Ha^{-n}$. Its measure is at most $K r^{-n}$ for a
constant $K$, and $r^{-n}=(a^{-n})^\alpha\le(a|I|/H)^\alpha$.
The finitely larger interval sizes are absorbed by increasing the
constant. Thus $\eta_{a,r}(I)\le C|I|^\alpha$, which proves the H\"older
bound for its distribution function.

For sharpness let $L$ be the left endpoint of the primitive support.
The leftmost level-$n$ cylinder lies in $[L,L+Ha^{-n}]$ and has exactly
mass $r^{-n}$; no other cylinder intersects this hull. Since there is no
atom at $L$, the distribution-function increment on this interval is
$r^{-n}$. The quotient of that increment by $(Ha^{-n})^\beta$ is
$H^{-\beta}a^{n\beta}/r^n$, which tends to infinity for $\beta>\alpha$.

For the smoothed density set $\lambda=\D_{a^{-d}}\eta_{a,r}$,
$\ell_j=a^{-j}$, and $R_d=\frac12\sum_{j<d}\ell_j$. Integrating
\eqref{eq:singular-derivative} once from $-\infty$ gives the explicit formula
\begin{equation}\label{eq:cdf-derivative}
 f^{(d-1)}(x)=\frac1{\prod_{j<d}\ell_j}
 \sum_{E\subseteq\{0,\ldots,d-1\}}(-1)^{|E|}
 F_\lambda\left(x+R_d-\sum_{j\in E}\ell_j\right).
\end{equation}
Each summand is H\"older of exponent $\alpha$, proving the upper assertion.
Near the left endpoint of the support of $f$, with distance smaller than
$\min_j\ell_j$, all terms except $E=\varnothing$ vanish. At that endpoint
the sharp increments just established for the scaled distribution function
therefore survive without cancellation, proving the lower assertion.
\end{proof}

The separated assumption is substantive. When $r\ge a$ and $a\nmid r$,
the measure is still singular by Theorem~\ref{thm:primitive}, but cylinder
intervals overlap and the bounded-intersection estimate used here is no
longer valid. No formula for its exact H\"older exponent is asserted.

\section{Verification, mathematical dependencies, and integration}
\label{sec:verification}

\subsection{What was checked computationally}
The accompanying \texttt{verify.py} uses only Python's standard library,
unbounded integers, and \texttt{fractions.Fraction}. Its exact finite tests
were run successfully for this article. The machine-readable receipt is
\texttt{verification\_results.json}. The principal test counts are:
\begin{center}
\begin{tabular}{@{}>{\raggedright\arraybackslash}p{10.1cm}r@{}}
\toprule
Test family & Cases or equalities\\
\midrule
Five-entry ladders with ratios in $\{2,3,4\}$, all ordered pairs,
scale denominators $1$ through $12$ & 78,732\\
Additional zero-order comparisons at fixed integer frequencies & 214,180\\
Exact uniform interval tilings & 96\\
Uniform splitting moments, through order $12$ & 1,248\\
Geometric finite-prefix factorization cases & 108\\
Finite-prefix factorization moments, through order $10$ & 1,188\\
Ordered pairs of subsets of an eight-element prime index set & 65,536\\
Allowed integer-base/scale triples checked at initial powers & 920\\
Forbidden integer-base/scale triples with explicit prime witnesses & 6,300\\
Composition-of-remainders moment equalities, through order $10$ & 825\\
Period-two ratio schedules and scale denominators $1$ through $24$ & 6,144\\
\bottomrule
\end{tabular}
\end{center}
The ladder tests compare coordinate divisibility, zero-order conditions at
source generators, and the lcm criterion. The moment tests check the
positive digit construction independently through exact rational moments.
For forbidden integer-base pairs the program computes a prime and an index
at which the zero-order inequality fails; it does not infer impossibility
from a finite absence of successful examples.

These are regression tests for algebra, indexing, normalization, and
finite instances. They are \emph{not} a replacement for the proofs of
infinite product convergence, singularity, Borel nonsmoothness, or hierarchy
completeness. No random sampling or floating-point experiment is used as
evidence for those conclusions.

\subsection{Proof dependencies and candidate formalization boundaries}
The mathematical dependencies can be separated into four layers. The first
is elementary integer arithmetic: divisibility chains, prime valuations,
initial-segment zero counts, and lcm obstructions. The second is probability:
independent bounded sums, finite uniform splitting, and uniqueness of
characteristic functions. The third is analysis: entire product convergence,
zero multiplicities, Fourier inversion, and finite-measure distributional
derivatives. The fourth is classification: prime coding, tail zero--one
arguments, and explicit computability reductions.

The repository's existing product and probability infrastructure is a
natural place to reuse established interfaces \cite{repo-inverse}. A sound
formalization plan should not simply postulate that a quotient is a
characteristic function. It should first formalize the uniform splitting
identity and assemble the actual independent digit law. Conversely, the
necessary direction should record compactness of the unknown factor before
using entire zero multiplicities.

A practical first formalization target is the arithmetic equivalence between
$\forall k\;(A_k\mid mB_k)$ and the generator zero-count inequalities.
A second target is the finite-prefix probability identity, followed by its
limit under an explicit deterministic tail bound. Those components already
isolate most of the main theorem's logical burden. The periodic-ratio
algorithm then reduces to finite arithmetic. The non-Borel classification
and hierarchy-completeness sections require additional libraries and should
not be labeled verified merely because the arithmetic component compiles.

\subsection{Result-status ledger}
\begin{center}
\begin{tabular}{@{}>{\raggedright\arraybackslash}p{5.2cm}>{\raggedright\arraybackslash}p{7.6cm}@{}}
\toprule
Result or ingredient & Status in this article\\
\midrule
Dyadic sinc product and dyadic zero orders & Classical; attributed to
Arias de Reyna.\\
Dyadic identical-convolution rootlessness & Identified repository overlap;
included only as a baseline corollary.\\
General ladder and cross-base factor classifications & Proved here from
explicit hypotheses; broader priority unconfirmed.\\
Prime coding in a compact smooth family; effective hardness & Proved here;
classical $E_0$ and computability methods are not claimed as new.\\
Primitive singularity and smoothing classification & Full proofs supplied;
related to classical self-similar-measure theory, with no exhaustive novelty
claim.\\
Sharp Fourier order and separated H\"older exponent & Proved here, including
noncancellation and resonance details.\\
Finite verification program & Executed successfully; not a Lean proof.\\
Further research questions & Proposed boundaries and extensions, not
assertions that their global open status has been fully audited.\\
\bottomrule
\end{tabular}
\end{center}

\section{Further research questions and topics}
\label{sec:questions}
The questions below specify extensions not settled by this article.
They should be checked against the relevant specialist literature before
being advertised as globally open problems. The aim is to identify concrete
next theorems and the precise mechanisms that the present proofs leave out.

\begin{question}[Beyond divisibility ladders]
For summable positive lengths $(\ell_k)$ and $(s_k)$, characterize when
$\prod_k\sinc(\ell_kt)/\prod_k\sinc(cs_kt)$ is the Fourier transform of a
probability measure. Which additional hypotheses make zero-divisor
domination sufficient, and can one construct a compactly supported
uniform-product counterexample where all zeros cancel but positivity fails?
\end{question}
Theorem~\ref{thm:ladder} uses an initial-segment property and coordinatewise
interval tilings. Without a ladder, a successful approach may require
matching several source factors to one target block, or a genuinely new
positive-factor criterion. A useful first target is a finite union of
integer divisibility chains with specified interlacing.

% ed. (2026-09-29): answered for prime-power geometric targets by a later draft.
\begin{ednote}
For a geometric target $\Mu{p^h}$ with $p$ prime and an arbitrary summable
family of uniform sources, the later unreviewed draft
\nolinkurl{Simultaneous_Convolution_Divisors_Fabius_Type_Laws/} of this
group (2026-09-29) proves that zero-divisor domination is sufficient: it is
equivalent to a nested matching of source factors to target scales and to
factorization by a probability measure. Its base-six theorem gives the
requested counterexample: in the notation of this article,
$\prod_{k\ge0}\sinc(t/6^k)\big/\bigl(\sinc(t/2)\sinc(t/3)\bigr)$ extends to
an entire function that is not the Fourier transform of a positive
measure, so $\Unif_{1/2}*\Unif_{1/3}$ does not divide $\Mu{6}$ although
all zeros cancel. That draft writes $X_p=\sum_k p^{-k}U_k$ with $U_k$
uniform on $[-1,1]$, so its $X_p$ is twice a variable with law $\Mu{p}$
here; it does not cite this article, whose one-ladder theorems are single
streams of its results. Targets that are not prime powers remain open.
\end{ednote}
% ed. (2026-09-30): a second sufficient hypothesis from a later draft (batch 69)
\begin{ednotelater}
The later unreviewed draft
\nolinkurl{Arithmetic_Rigidity_off_Resonance_Geometric_Uniform_Laws/} of
this group (batch~69 of the repository's \nolinkurl{docs/incoming/},
2026-09-30), written in the normalization of this article, gives a second
hypothesis under which zero-divisor domination is sufficient. If the target
lengths $\ell_k$ are pairwise rationally incommensurable, a finite or
countable summable family of uniforms divides the law of
$\sum_k\ell_kU_k$ exactly when the Fourier quotient is entire, and exactly
when the uniforms have lengths $\ell_{\kappa(j)}/n_j$ for an injective
assignment $\kappa$ and integers $n_j\ge1$ (its \texttt{thm:separated}).
Its \texttt{thm:channels} gives the same conclusion for the geometric
targets $\mu_p$ with $p=r^{-e/h}$, $r$ prime, through $h$ independent
prime-adic Hall conditions.
\end{ednotelater}

\begin{question}[Two arbitrary geometric ratios]
Determine all triples $(p,q,c)$ with $0<p,q<1$, $c>0$, for which
$\mu_p=\D_c\mu_q*\nu$ has a probability solution.
\end{question}
Theorem~\ref{thm:self-spectrum} treats $p=q$, and
Theorem~\ref{thm:crossbase} treats $p=1/a$ for an integer $a$.
For general $p$, the target zeros are no longer confined to integers.
The necessary inclusion of multiplicative zero sets suggests a role for
relations between powers of $p$ and $q$, but zero inclusion alone has not
been proved sufficient in that setting.
% ed. (2026-09-30): answered for all but countably many target ratios by a later draft (batch 69)
\begin{ednotelater}
The later unreviewed draft
\nolinkurl{Arithmetic_Rigidity_off_Resonance_Geometric_Uniform_Laws/} of
this group (batch~69, 2026-09-30; same normalization) answers this
question for every target ratio $p$ no positive power of which is
rational, a co-countable set: $\mu_p=\D_c\mu_q*\nu$ has a probability
solution exactly when $c=p^a/n$ and $q=p^m/d$ with integers $a\ge0$ and
$m,n,d\ge1$ (its \texttt{thm:geometric}). For $p=r^{-e/h}$ with $r$ prime
and $\gcd(e,h)=1$, its \texttt{thm:orbit} gives a finite criterion: the
least $\ell\ge1$ with $q^\ell$ rational exists and divides $h$, the widths
$cq^j$, $0\le j<\ell$, lie in distinct rationality classes
$p^{s}\mathbb Q$, $0\le s<h$, and $q^\ell=1/D$ for an integer $D$ divisible by
$r^e$. In both cases zero inclusion is sufficient. Together with
Theorem~\ref{thm:crossbase}, the question remains open only for targets
whose least rational power $p^h$ is $A/B$ in lowest terms with $A>1$, or
$1/B$ with $h\ge2$ and $B$ not a prime power. That draft re-proves
Theorem~\ref{thm:self-spectrum} as its \texttt{prop:same}, without citing
this article.
\end{ednotelater}

\begin{question}[Regularity of cross-base digit remainders]
When $a\mid b$ and $b>a$, classify the absolute continuity, exact smoothness,
and Fourier decay of $\nu_{(a^k),(b^k),m}$.
\end{question}
Here digit counts grow as $m(b/a)^k$, so the remainder is not the
fixed-digit self-similar law treated in Sections~\ref{sec:regularity}
and~\ref{sec:fourier}. The finite digit mesh becomes much smaller than the
width it fills. A quantitative comparison with an infinite uniform product
might yield smoothing beyond the same-base classification, but this
requires a proof controlling all frequency scales.

\begin{question}[Overlapping primitive digits]
For $r\ge a$ with $a\nmid r$, find effective formulas or algorithms for
local dimensions and optimal H\"older exponents of the distribution
function of $\eta_{a,r}$.
\end{question}
The measure has interval support and is singular; neither fact determines
its local exponents. Finite carry-state descriptions of base-$a$
expansions offer one concrete route. The separated-case exponent
$\log r/\log a$ cannot simply be reused: it is at least one in this range
and the separated-cylinder proof breaks down.

\begin{question}[Finer function-space regularity]
Determine sharp Sobolev, Besov, and fractional-variation membership for the
densities of $\eta_{a,a^dr}$, especially when the primitive digits overlap.
\end{question}
The article fixes the last continuous integer derivative and proves a
sharp pointwise power Fourier envelope. These do not alone determine
integrated Fourier norms or all fractional smoothness scales, because a
resonant subsequence can be thin. Quantifying the width and frequency of
resonance neighborhoods is a possible next step.

\begin{question}[Optimal approximate factorization]
Determine or sharply estimate
$\inf_\nu W_1(\mu_A,\D_{1/m}\mu_B*\nu)$, and identify when the Fourier
obstruction in Proposition~\ref{prop:wasserstein} is asymptotically optimal.
\end{question}
The proposition already gives a strictly positive, support-independent
lower bound from any failed divisibility condition. Remaining targets
include upper bounds from constructive approximate factors, optimization
over witness frequencies, and analogous sharp estimates in other metrics.
Any estimate must depend on the witness scale:
Proposition~\ref{prop:continuous} shows how small changes in high prime
positions can be in every fixed derivative norm. A uniform separation over
all nonfactorable pairs is therefore not a plausible substitute for a
scale-sensitive analysis.
% ed. (2026-09-30): related quantitative results in a later draft (batch 69)
\begin{ednotelater}
For geometric targets off resonance, the batch-69 draft named in the
notes after ``Beyond divisibility ladders'' and ``Two arbitrary geometric
ratios'' bounds $\inf_\nu W_1(\mu_q,\Unif_{1/N}*\Unif_{1/M}*\nu)$ from
below by the mechanism of Proposition~\ref{prop:wasserstein} (its
\texttt{thm:metric}, which does not cite this article). For
$\Unif_{1/2}*\Unif_{1/3}$ it shows that the infimum is $0$ at $q=1/2$,
positive at every $q$ no positive power of which is rational, and at most
$|q-1/2|/(2(1-q))$ (its \texttt{thm:perturbation}), so exact factorability
is lost under arbitrarily small changes of $q$ while the distance tends to
zero. The sharp rate is open there too.
\end{ednotelater}

\begin{question}[Decidable presentations beyond eventual periodicity]
Extend Proposition~\ref{prop:periodic} to ratio sequences given by finite
automata, substitutions, or other certified finite-state presentations.
\end{question}
The arithmetic problem is to decide whether all prime-exponent deficits
remain bounded above and whether only finitely many primes contribute.
For a finite alphabet of ratios, the second issue disappears, leaving a
finite collection of weighted prefix-sum boundedness problems. The precise
answer depends on the presentation class and should be stated with its
input model and complexity, rather than as a generic algorithm for all
computable ladders.

\begin{question}[Classification complexity on larger classes]
What is the exact Borel classification complexity of mutual scaled
convolution factorability for arbitrary smooth compactly supported
probability laws, or for broader infinite uniform products?
\end{question}
Theorem~\ref{thm:nonsmooth} supplies a concrete $E_0$ lower bound, even on a
compact family. It does not establish a maximal lower bound or an upper
classification theorem for the ambient class. A useful intermediate
problem is to determine the complexity for finitely many interlaced ladders.

\begin{question}[Multivariate arithmetic factors]
Develop an analogue of Theorem~\ref{thm:ladder} for sums of uniforms on
nested boxes, parallelotopes, or lattice fundamental domains in $\R^d$.
\end{question}
For products of independent one-dimensional ladders, the theorem applies
coordinatewise when the dilation is diagonal. Non-diagonal integer maps
introduce genuine lattice geometry. Smith normal forms may describe finite
uniform tilings, but the Fourier zero sets are hyperplanes rather than
isolated points, and multiplicity bookkeeping must be redesigned.

% ed. (2026-09-29): product-target case treated by a later draft.
\begin{ednote}
For the independent product of prime-base geometric targets
${\Mu{p}}^{\otimes d}$, the later unreviewed draft
\nolinkurl{Simultaneous_Convolution_Divisors_Fabius_Type_Laws/} of this
group (2026-09-29) proves that a sum of uniform laws along arbitrary
vectors in $\R^d$ is a convolution factor exactly when every generator is
coordinate-aligned with a signed reciprocal-integer length and each
coordinate satisfies the one-dimensional matching criterion, and that the
real matrices $C$ for which $C_*{\Mu{p}}^{\otimes d}$ is a convolution factor
of ${\Mu{p}}^{\otimes d}$ are exactly the partial monomial matrices with
nonzero entries $\pm1/n$. Nested boxes, parallelotopes and lattice
fundamental domains in general remain open.
\end{ednote}

\begin{question}[Certified probability-factor software]
Implement an exact factorization engine that takes a certified ladder
presentation, returns either an explicit zero-order obstruction or a
proof of the infinite divisibility condition, and produces a verified
sampler or weak approximation for the remainder.
\end{question}
The unrestricted decision problem is impossible by
Theorem~\ref{thm:complexity}, so the specification must permit an
``undetermined'' result outside certified subclasses. Positive outputs
should include a constructive digit law, not only cancellation of an
entire quotient. The finite-prefix test suite in this package is a
regression starting point, while the convergence and positivity statements
are the essential proof obligations.

\section{Conclusion}
The classification problem admits an unusually exact answer once its
arithmetic structure is exposed. For nested integer denominators,
probability factorization is equivalent to divisibility at every level;
the remainder is then an explicit independent-digit sum. The same criterion
simultaneously yields elementary cross-base rigidity, a finite algorithm
for periodic schedules, and sharp nonclassification results for unrestricted
presentations.

This coexistence is the main conceptual point. The targets are smooth,
compactly supported, and effectively approximable, yet their exact
convolution relation can remember infinitely many independent arithmetic
choices. The remainder measures also retain that structure analytically:
removing a scaled smooth factor can leave a singular law, a finite spline,
or a density whose highest continuous derivative has singular variation.
The written proofs and finite exact checks establish these statements
within the hypotheses given here; the broader questions and priority
assessment remain separate research tasks.

\appendix
\section{Reproduction and source notes}
\label{sec:reproduction}
The package is self-contained for typesetting. From its directory, run
\begin{verbatim}
python3 verify.py --output verification_results.json
pdflatex -interaction=nonstopmode -halt-on-error article.tex
pdflatex -interaction=nonstopmode -halt-on-error article.tex
pdflatex -interaction=nonstopmode -halt-on-error article.tex
\end{verbatim}
The three TeX passes resolve the contents and cross-references. A
\texttt{Makefile} is included as a convenience. The PDF supplied with this
package was generated from this TeX source, not copied from the repository.
No external figures, downloaded font files, proprietary packages, or network
requests are required. Standard LaTeX packages listed in the preamble are
needed.

\texttt{SOURCES.md} records the inspected repository locations and primary
bibliographic sources. The repository snapshot is identified by the commit
printed in Section~\ref{sec:introduction}; subsequent changes are outside
this audit. The standalone article is an addition for consideration, not
a commit, pull request, or modification of the user's repository.

\begingroup
\small
\setlength{\parskip}{0pt}
\begin{thebibliography}{9}
\bibitem{fabius}
J. Fabius,
\emph{A probabilistic example of a nowhere analytic $C^\infty$-function},
Zeitschrift f\"ur Wahrscheinlichkeitstheorie und verwandte Gebiete
\textbf{5} (1966), 173--174.
CWI archival record and report: \url{https://ir.cwi.nl/pub/8160}.

\bibitem{arias}
J. Arias de Reyna,
\emph{An infinitely differentiable function with compact support:
Definition and properties},
English translation (2017) of the 1982 article in
Revista de la Real Academia de Ciencias Exactas, F\'isicas y Naturales
\textbf{76}, 21--38.
\url{https://arxiv.org/abs/1702.05442}.
In particular, equations (4) and (9) give the dyadic Fourier product and
integer-zero multiplicities.

\bibitem{hutchinson}
J. E. Hutchinson,
\emph{Fractals and self-similarity},
Indiana University Mathematics Journal \textbf{30} (1981), no.~5, 713--747.
\url{https://doi.org/10.1512/iumj.1981.30.30055}.

\bibitem{hkl}
L. A. Harrington, A. S. Kechris, and A. Louveau,
\emph{A Glimm--Effros dichotomy for Borel equivalence relations},
Journal of the American Mathematical Society \textbf{3} (1990), no.~4,
903--928.
\url{https://doi.org/10.1090/S0894-0347-1990-1057041-5}.
Author-repository record:
\url{https://authors.library.caltech.edu/records/cvb33-qt212}.

\bibitem{repo-inverse}
V. Reshetnikov and contributors,
\emph{ProveIt: Inverse and sampling}, Fabius-function research documentation,
inspected at commit \sourcecommit.
\href{https://github.com/VladimirReshetnikov/ProveIt/blob/fbba58593dc0622aa914972896150d4848f935b5/Analysis/FabiusFunction/docs/semi-formalized-research-frontiers/drafts/inverse-and-sampling/README.md}{Pinned repository README}.

\bibitem{repo-shape}
V. Reshetnikov and contributors,
\emph{ProveIt: Shape, divisibility, and Stein geometry of the Fabius--Rvachev law},
research-package README, inspected at the same commit.
\href{https://github.com/VladimirReshetnikov/ProveIt/blob/fbba58593dc0622aa914972896150d4848f935b5/Analysis/FabiusFunction/docs/semi-formalized-research-frontiers/drafts/representations/Fabius_Rvachev_Shape_Divisibility_Stein_Geometry/README.md}{Pinned repository README}.
This source explicitly identifies the identical-convolution rootlessness
strand and distinguishes manuscript proofs from Lean verification.
\end{thebibliography}
\endgroup
\end{document}
